\beginsong{Love}[by={Jean-Claude Gianadda}]
\textnote{Capo 5}
\beginverse
\[Am]Il est venu chez nous, guitare en bandou\[Dm]lière,
Venait d'on ne sait \[G]où, il parcourait la \[C]terre.
Et dans ses longs \[E]cheveux, le vent semblait chan\[Am]ter,
Tout au fond de ses \[E]yeux dansait la liber\[Am]té.
\endverse
\beginchorus
\[Am]Love, c'était son \[Dm]nom,(lalala) \[G]Love,
Un vaga\[C]bond qui vivait de so\[E]leil,
D'espace et de chan\[Am]sons.
\endchorus
\beginverse
Il écoutait le \[Am]vent, les fleurs et les ri\[Dm]vières,
Jouait comme un en\[G]fant, parlait de la lu\[C]mière.
Il partageait ses \[E]rires, ses rêves et ses pro\[Am]jets,
Et dans chaque sou\[E]rire dansait la liber\[Am]té.
\endverse
\beginverse
Il est parti un \[Am]jour, nul ne sait où il \[Dm]est,
Au pays de l'a\[G]mour, tu peux le rencon\[C]trer.
Et dans notre mai\[E]son, il nous aura lais\[Am]sé,
Avec cette chan\[E]son, un peu de liber\[Am]té.
\endverse
\endsong
